\subsubsection{Partícula atada a una cuerda sobre un plano inclinado}

Una partícula $A$ de masa $m$ está unida a un extremo de una cuerda ideal de longitud $R$,
cuyo otro extremo se fija a un punto del plano. La partícula describe una circunferencia de
radio $R$ contenida en un plano inclinado un ángulo $\phi$ respecto a la horizontal, como se
muestra en la Figura \ref{fig:c2_e01_p20_cuna}. El ángulo $\theta$ mide la posición angular
de la partícula sobre la circunferencia, medido desde el punto más alto $B$. Existe un
coeficiente de fricción deslizante $\mu_c$ entre $A$ y la superficie inclinada, la gravedad
actúa verticalmente hacia abajo, y la partícula pasa por el punto $B$ con rapidez $v_0$
tangente a la circunferencia.

\begin{enumerate}
    \item Defina marcos de referencia y sistemas de coordenadas apropiados.
    \item Determine la ecuación diferencial que gobierna el ángulo $\theta$.
    \item Calcule la rapidez $v_0$ mínima en $B$ necesaria para dar exactamente una vuelta
          completa.
\end{enumerate}

\begin{figure}[!ht]
    \centering
    \begin{tikzpicture}[>=Latex, scale=1.1]

        % --- VISTA EN PLANTA (plano inclinado en perspectiva) ---
        % --- PARAMETROS ---
        \def\Rx{2.6}   % semieje mayor (perspectiva)
        \def\Ry{1.1}   % semieje menor (perspectiva)
        \def\ang{125}  % angulo theta de A

        % Contorno del plano inclinado (paralelogramo)
        \draw[thick] (-3.4,-1.7) -- (3.4,-1.7) -- (4.2,1.7) -- (-2.6,1.7) -- cycle;

        % Circunferencia (trayectoria) en perspectiva
        \draw[dashed] (0,0) ellipse ({\Rx} and {\Ry});

        % Centro y punto B (mas alto)
        \coordinate (O) at (0,0);
        \coordinate (B) at (0,\Ry);
        \coordinate (A) at ({\Rx*cos(\ang)},{\Ry*sin(\ang)});
        \filldraw (O) circle (1pt);
        \filldraw (B) circle (1.5pt) node[above] {$B$};
        \filldraw (A) circle (1.5pt) node[left] {$A\;(m)$};

        % Cuerda y referencia
        \draw[thick] (O) -- (A) node[midway, below] {$R$};
        \draw[dashed] (O) -- (B);

        % Angulo theta
        \draw[->] (0,0.7) arc (90:\ang:0.7);
        \node at (155:1.05) {$\theta$};

        % --- VISTA LATERAL (cuña) ---
        \begin{scope}[shift={(7.2,-1.2)}]
            \draw[thick] (0,0) -- (3.2,0) -- (3.2,1.7) -- cycle;
            \draw[->] (0.55,0) arc (0:28:0.55);
            \node at (0.95,0.18) {$\phi$};
            \filldraw (1.5,0.8) circle (1.2pt);
            \filldraw (2.1,1.12) circle (1.2pt) node[above left=-1pt] {$B$};
            \node at (1.55,1.15) {$R$};
            \draw[->, thick] (-0.4,1.6) -- (-0.4,0.8) node[midway,left] {$g$};
            \node[font=\small] at (1.6,-0.45) {vista lateral};
        \end{scope}

    \end{tikzpicture}
    \caption{Partícula $A$ que recorre una circunferencia de radio $R$ sobre un plano inclinado un ángulo $\phi$; vista en planta y vista lateral.}
    \label{fig:c2_e01_p20_cuna}
\end{figure}
